\section{Introduction}\label{sec:intro}
A finite correspondence between two group factors is a finite-dimensional
relation in the sense of von Neumann module dimension. For lattice groups,
we prove that this analytic relation comes from a finite collection of
actual subgroup isomorphisms. Scalar cocycles remain visible: the finite
fibres are projective representations whose multiplier is exactly the
ratio of the two cocycles along the subgroup isomorphism.

The source groups include lattices in products of higher-rank real and
nonarchimedean groups, together with the quaternionic and Cayley
property-$(T)$ rank-one groups. The comparison group is only assumed to
be countably infinite and ICC. Thus the geometry used in the proof must
be constructed inside the correspondence; no boundary or ambient Lie
group is initially available on the comparison side.

\subsection{The class, conventions, and main theorem}
\begin{definition}\label{def:class}
Let $\mathscr K$ be the class of countably infinite ICC groups abstractly
commensurable with a lattice in a finite nonempty product
\[
 G=\prod_{i=1}^r G_i,\qquad G_i=\mathbf H_i(k_i)^+.
\]
Here $k_i$ is $\mathbb R$, $\mathbb C$, or a finite extension of
$\mathbb Q_p$; the primes may differ. The algebraic group $\mathbf H_i$ is
connected, adjoint, and absolutely simple over $k_i$. The superscript $+$
denotes the subgroup generated by the $k_i$-points of the unipotent radicals
of proper $k_i$-parabolic subgroups. Each $G_i$ is required to be noncompact
and to have Kazhdan's property $(T)$.

A lattice is discrete and has finite Haar covolume. Abstract
commensurability means an actual isomorphism between finite-index
subgroups. ICC means that every nonidentity conjugacy class is infinite.
No irreducibility, cocompactness, product decomposition, or torsion
hypothesis is imposed on the lattice.
\end{definition}

Property $(T)$ means that every continuous unitary representation with
almost invariant unit vectors has a nonzero invariant vector. The real
simple factors in Definition~\ref{def:class} have higher rank or are locally
$\Sp(n,1)$, $n\geq2$, or $F_{4(-20)}$. At a finite place their relative
rank is at least two. Complex places are treated as real Lie groups when
using differential arguments. Standard structure and permanence results
are recalled in Section~\ref{sec:setup}; see also \cite{BHV2008}.

All scalar cocycles take values in $\mathbb T$ and are normalized:
\[
 \mu(g,h)\mu(gh,j)=\mu(h,j)\mu(g,hj),\qquad
 \mu(e,g)=\mu(g,e)=1.
\]
The twisted group factor $L_\mu(\Gamma)$ has canonical unitaries $u_g$
with $u_gu_h=\mu(g,h)u_{gh}$ and trace
$\tau(u_g)=\mathbf1_{\{e\}}(g)$. For an ICC group it is a factor: the
absolute values of Fourier coefficients of a central vector are constant
on conjugacy classes, so square summability kills all nonidentity
coefficients.

A bifinite $M$--$N$ correspondence is a Hilbert space with commuting normal
unital left $M$ and right $N$ actions and finite von Neumann module
dimensions on both sides. The actions need not be mutual commutants, and
the correspondence need not be irreducible. Every correspondence in this
paper is separable; every factor isomorphism is normal, unital,
complex-linear, and $*$-preserving. For a finite factor $Q$, write
\[
 Q^t=p\bigl(M_n(\mathbb C)\,\overline\otimes\,Q\bigr)p,
 \qquad (\Tr_n\otimes\tau_Q)(p)=t>0.
\]
The matrix trace $\Tr_n$ is unnormalized; the corner itself has its
normalized trace.


For finite-index subgroups $A\leq\Gamma$ and $B\leq\Lambda$, an
isomorphism $\delta:A\to B$, and a nonzero finite-dimensional projective
unitary representation $\sigma$ satisfying
\begin{equation}\label{eq:intro-multiplier}
 \sigma(a)\sigma(b)
 =\frac{\mu(a,b)}{\omega(\delta(a),\delta(b))}\sigma(ab),
 \qquad \sigma(e)=I,
\end{equation}
define
\begin{equation}\label{eq:intro-model}
 E(A,B,\delta,\sigma)
 =L^2(M)\,\overline\otimes_{L_{\mu|_A}(A)}
       \bigl(V_\sigma\otimes L^2(N)\bigr),
 \qquad M=L_\mu(\Gamma),\quad N=L_\omega(\Lambda).
\end{equation}
The right $N$ action is standard. The left subgroup action on the second
factor is
$u_a(\xi\otimes\eta)=\sigma(a)\xi\otimes v_{\delta(a)}\eta$.
Equation~\eqref{eq:intro-multiplier} makes this an action with multiplier
$\mu|_A$.


A right action is represented by $\rho(h)\xi=\xi v_h^*$, so
$\rho(h)\rho(k)=\overline{\omega(h,k)}\rho(hk)$.
The left action is $\pi(g)\xi=u_g\xi$.  Thus the joint action
$U_{(g,h)}\xi=u_g\xi v_h^*$ has multiplier
$\alpha((g,h),(g',h'))=\mu(g,g')\overline{\omega(h,h')}$.

\begin{theorem}[Arithmetic exhaustion]\label{thm:arithmetic-exhaustion}
Let $\Gamma\in\mathscr K$, let $\Lambda$ be any countably infinite ICC
group, and let $\mu\in Z^2(\Gamma,\mathbb T)$ and
$\omega\in Z^2(\Lambda,\mathbb T)$ be normalized scalar cocycles.
Every bifinite $L_\mu(\Gamma)$--$L_\omega(\Lambda)$ correspondence is
unitarily isomorphic to a closed bimodule summand of a finite direct sum
of the models $E(A,B,\delta,\sigma)$ in Equation~\eqref{eq:intro-model},
where $A\leq\Gamma$ and $B\leq\Lambda$ have finite index,
$\delta:A\to B$ is an actual group isomorphism, and $\sigma$ is a
finite-dimensional projective unitary representation satisfying
Equation~\eqref{eq:intro-multiplier}.

Conversely, every such model, finite direct sum, and closed bimodule
summand is a normal bifinite correspondence. If
$n=\dim_{\mathbb C}V_\sigma$, the elementary model has dimensions
\[
 \dim_{L_\mu(\Gamma)}E(A,B,\delta,\sigma)=n[\Lambda:B],\qquad
 \dim E(A,B,\delta,\sigma)_{L_\omega(\Lambda)}=n[\Gamma:A].
\]
\end{theorem}

The theorem retains both finite sums and closed summands. Its displayed
dimension formulas concern the elementary models themselves. The finer
problem of identifying every summand and every bounded intertwiner is
addressed in the companion article \cite{OpenAI2026Category}, which
uses Theorem~\ref{thm:arithmetic-exhaustion} as its geometric input.
The elementary construction and its normality and dimensions are proved
here in Section~\ref{sec:models}.

Only the matched ratio $\mu/\delta^*\omega$ must admit the
finite-dimensional projective representation in
Equation~\eqref{eq:intro-multiplier}. Neither endpoint cocycle is required
to have finite order, finite type, or a virtually trivial restriction,
and the image of $\sigma$ may be infinite. In particular, existence of a
nonzero correspondence forces $\Gamma$ and $\Lambda$ to be abstractly
commensurable. For trivial cocycles the converse follows by taking the
trivial one-dimensional fibre. These conclusions apply to nonuniform
and reducible lattices as well as to their ICC commensurable groups.

\subsection{Rigidity and finite-correspondence precedents}
The group-factor reconstruction problem asks which features of a
discrete group survive in its von Neumann algebra. Connes proved
countability of the fundamental group of a property-$(T)$ group factor
\cite{Connes1980}; his later problems ask whether nonisomorphic ICC
property-$(T)$ groups have nonisomorphic factors and ask for the
fundamental groups of those factors
\cite[Chapter~V, Appendix~B.$\varepsilon$, Problems~1--2]{Connes1994}.
Popa's survey states a stronger recovery question that controls the given
isomorphism, up to a character and inner conjugacy, together with its
amplification scale
\cite[Section~3, Equations~(3.2) and~(3.2$'$)]{Popa2007}.
Property $(T)$ alone does not ensure group reconstruction: independent,
concurrent constructions of OpenAI and Zhou give nonisomorphic ICC
property-$(T)$ groups with isomorphic group factors
\cite[Theorem~1.2]{OpenAI2026ConnesCounterexample}
\cite[Theorem~A]{Zhou2026Counterexamples}.
The present theorem concerns the lattice class in
Definition~\ref{def:class}. Arithmetic exhaustion is the
finite-correspondence input used in the companion article to obtain
stable canonical recovery for this class.

Deformation and rigidity methods have produced strong reconstruction
theorems in several settings. Ioana's theorem recovers Bernoulli
actions of ICC property-$(T)$ groups from their crossed products
\cite{Ioana2011}. Ioana--Popa--Vaes obtained $W^*$-superrigid groups
among generalized wreath products \cite{IPV2013}, and
Chifan--Ioana--Osin--Sun constructed property-$(T)$ examples using
wreath-like products \cite{CIOS2023}. These results recover groups or
actions from operator-algebraic data through constructions different
from the lattice geometry developed here.

Finite correspondences retain finite-index relations and their
compositions. Jones' index and module dimensions \cite{Jones1983},
the finite-index expectation inequality of Pimsner--Popa
\cite{PimsnerPopa1986}, and Connes' relative tensor product
\cite[Chapter~V, Appendix~B]{Connes1994} give the framework.
Vaes' classification for generalized Bernoulli factors already organizes
finite-index bimodules by subgroup isomorphisms and finite-dimensional
projective representations of cocycle ratios, and computes their fusion
by projective induction \cite[Theorem~2.2 and Section~6.2]{Vaes2008}.
Donvil--Vaes prove rigidity under arbitrary scalar twists and virtual
isomorphism for certain left-right wreath products. Their results
include the projective-ratio constraint and reconstruction, up to
amplification, of finite-index factor neighbors
\cite[Theorems~A and~6.6]{DonvilVaes2025}.
Chifan--Fern\'andez Quero--Osin--Tan prove stable canonical recovery
for twists of certain property-$(T)$ wreath-like products when the
source cocycle is trivial on the abelian base
\cite[Theorem~F]{CFOT2026}. The exhaustion theorem here treats arbitrary
source cocycles in the local-field lattice class.

A geometric line of operator-algebraic rigidity for lattices begins
with character and finite-factor representation results. Bekka
classified the characters of $\SL_n(\mathbb Z)$, $n\geq3$
\cite[Theorem~3]{Bekka2007}; Peterson established a finite-factor
representation dichotomy for irreducible higher-rank property-$(T)$
lattices \cite[Theorem~A]{Peterson2014}.
Boutonnet--Houdayer developed stationary and noncommutative boundary
methods that recover character rigidity
\cite[Theorems~A--C]{BoutonnetHoudayer2021}.
Their noncommutative factor theorem subsequently classified
intermediate von Neumann algebras of specified lattice boundary
crossed-product inclusions by parabolic quotients
\cite[Theorem~B and Theorem~3.3]{BoutonnetHoudayer2023}.
Related results recover Weyl groups from the distinguished inclusion
associated with the action on the quotient by a split-torus centralizer
\cite[Theorem~A and Corollary~B]{HI2026}, and classify intermediate
$C^*$-algebras for boundary crossed products of irreducible lattices in
products of real rank-one groups
\cite[Corollary~B]{HoudayerLeBars2026}.
These results show how boundary data can retain the ambient geometry.
For the finite correspondence considered here, compatible boundary
measurements must be constructed in several module copies; their
simultaneous law then produces finite-rank subgroup-graph fibres.

The search for a classical equivariant map from an operator-valued
coupling is motivated in part by the self-coupling strategy in Furman's
measure-equivalence rigidity theorem \cite[Sections~3--4]{Furman1999}.
The noncommutative classicalization and finite-module reconstruction
used here are proved below; they are not consequences of that classical
theorem.

\subsection{From boundary measurements to subgroup graphs}
The proof first passes to an actual lattice and removes the compact
projective sectors of the finite commutant. The remaining conjugation
action is weakly mixing. In this reduced setting the goal is a fixed
orthogonal partition of the correspondence $K$ into finite-dimensional
subspaces, permuted by both endpoint groups, with free actions and
finitely many orbits. A joint stabilizer of one subspace is then the graph
of an actual finite-index subgroup isomorphism, and its action on that
subspace is the projective fibre in
Equation~\eqref{eq:intro-multiplier}.

The geometric construction produces this partition through boundary
measurements. Write $B$ for the compact full flag space of the known
ambient group $G$. A projection-valued measure, or PVM, assigns an
orthogonal projection to each Borel set, with countable additivity in
the strong operator topology; disjoint sets give orthogonal projections
and the whole space gives the identity. An angular measurement $P_x$
on $B$ therefore records the possible image of a source flag $x\in B$:
$P_x(A)$ is the projection for its image to lie in $A\subset B$.
Initially the construction supplies positive operator-valued
observations. Proving that they are PVMs, and that observations at
different labels commute, are separate parts of the argument.

To compare source copies, use $K_0\otimes K_1$, with the left
$\Gamma$-action in each slot and the right $\Lambda$-action acting
diagonally. The eventual joint angular law is a PVM $\Theta$ on the
space of nonsingular, essentially invertible maps $B\to B$. The two
left actions transport its source and target labels; diagonal
$\rho(\Lambda)$ fixes it. The same construction on
$K_0\otimes K_1\otimes K_2$ compares slot zero separately with slots
one and two. The required conclusion includes
\[
 [\Theta_{01}(A),\Theta_{02}(C)]=0
 \qquad\text{for all Borel sets of boundary maps }A,C.
\]
This common-slot commutation has a concrete purpose. Normal slices in
the separate branch slots generate an abelian algebra on $K_0$.
Disintegrating over it gives boundary-map PVMs on the fixed space $K_1$.
Finite module dimension makes their orbit transversals finite rank;
weak mixing makes the resulting projection partition independent of
the disintegration parameter. The boundary-map labels themselves need
not become deterministic. Section~\ref{sec:recovery} proves this passage
and the subsequent graph-stabilizer identification.

Both copies of $B$ in this construction belong to the source lattice.
No flag space for the unknown group $\Lambda$ is assumed. To read the
known geometry from Fourier coefficients, omit one left action and put
the remaining actions, together with diagonal $\rho$, in a finite
regular column. Inducing the omitted action to $G/\Gamma$ gives a
finite tracial algebra even when the lattice is reducible. Escape in
all ambient factors then creates an isometric boundary link. Escape
in a subproduct instead gives a pairing through the fixed algebra's
conditional expectation. Ordered Fourier transfer and incidence
relations make the positive observations sharp and identify exact
residue graphs, while unrelated measurements still retain their
original order.

Three local mechanisms control the missing Cartan gaps. At finite
places, compact-open fixed algebras and a uniform subgroup logarithm
argument recover homogeneous residue maps and upper speed bounds.
Real panels inside higher-rank components use optimal finite-dimensional
modules: trace inequalities force a nonzero scalar witness in their
section spaces, and positivity of its algebraic action excludes excess
speed. Isolated quaternionic and Cayley factors use uniform ray
estimates and a finite-module spectral gap to obtain positive mass in
linear windows on fixed vectors. Pair packing then gives horizontal
derivatives, whose closed module core produces an injective positive
height with its own-label covariance. The Carnot boundary geometry
belongs to the work of Cowling--Dooley--Kor\'anyi--Ricci and Pansu
\cite{CDKR1991,CDKR1998,Pansu1989}; the ray and finite-module estimates
used here are proved in Sections~\ref{sec:ray}--\ref{sec:scale}.

Full tuples bring these branches together. Comparing lattice packing
with Haar growth calibrates the speeds around the finite type
permutation and fixes the exact Jacobians. Ordered imaginary powers
then produce a Haar-height isometry. Its common graph core permits
locality arguments for the unbounded heights. Locality gives finite
jets for the actual spectral control measure, which may be singular;
commuting cone multipliers eliminate every positive-order symbol.
The principle is related to Peetre's locality theorem
\cite{Peetre1959,Peetre1960}, but the measured operator-domain argument
is proved here. Applying it also to the two shared-slot copies gives
the isolated-type commutation needed for the joint law.

For higher-component panels, pairing powers and exact residue graphs
supply the corresponding commutation. Simultaneous reflections and
panel galleries then glue all local types into the nonsingular,
essentially invertible chamber maps described above. After extracting
the fixed finite-rank partition and its graph fibres, we restore every
compact sector with its exact projective multiplier. Two-sided
finite-index induction returns the original correspondence as a closed
summand. This explains both the matched cocycle ratio and the summand
qualification in Theorem~\ref{thm:arithmetic-exhaustion}.

\begin{figure}[tbp]
 \centering
 \begin{tikzpicture}[
  box/.style={draw=linkblue!65,rounded corners=2pt,fill=linkblue!4,
    align=center,text width=40mm,minimum height=16mm,inner sep=4pt,
    font=\small},
  wide/.style={box,text width=137mm},
  arrow/.style={-{Stealth[length=2mm]},draw=linkblue!85,line width=.6pt},
  node distance=5mm and 4mm]
\node[wide] (links) {Finite modules and ordered boundary links
  (Sections~\ref{sec:setup}--\ref{sec:transfer})\\
  Remove compact projective sectors; construct complementary regular models.\\
  Count gaps and obtain sharp flags with exact residue graph relations.};
\node[box,below=of links] (real) {Real panels in\\higher-rank factors\\
  Section~\ref{sec:panels}\\
  Finite-dimensional stages\\and upper gap bounds};
\node[box,left=of real] (finite) {Finite-place panels\\
  Section~\ref{sec:finite}\\
  Compact-open algebras,\\homogeneous pencils,\\and upper gap bounds};
\node[box,right=of real] (isolated) {Isolated real factors\\
  Sections~\ref{sec:ray}--\ref{sec:scale}\\
  Ray densities, linear\\subsequences, and heights\\at each vertex};
\node[wide,below=of real] (haar) {Product calibration and Haar-height isometry
  (Section~\ref{sec:heights})\\
  Compare full-tuple Haar densities; obtain tight centered gaps,\\
  inverse-height integrability, and exact height intertwining.};
\node[wide,below=of haar] (locality) {Isolated-type locality
  (Section~\ref{sec:locality})\\
  Construct dense graph cores for the actual spectral measures;\\
  eliminate positive-order jets and prove commutation in both normal copies.};
\node[wide,below=of locality] (gluing) {Pencils and simultaneous chamber laws
  (Section~\ref{sec:gluing})\\
  Establish real-pencil commutation and glue all local types.\\
  Obtain jointly commuting, nonsingular, essentially invertible boundary maps.};
\node[wide,below=of gluing] (arithmetic) {Arithmetic exhaustion
  (Section~\ref{sec:recovery})\\
  Uninduce the joint law; recover deterministic finite-rank atoms and graph fibres.\\
  Restore compact sectors and finite-index descent: finite sums and closed summands.};
\draw[arrow] (links.south) -- (real.north);
\draw[arrow] (links.south -| finite) -- (finite.north);
\draw[arrow] (links.south -| isolated) -- (isolated.north);
\draw[arrow] (real.south) -- (haar.north);
\draw[arrow] (finite.south) -- (haar.north -| finite);
\draw[arrow] (isolated.south) -- (haar.north -| isolated);
\draw[arrow] (haar) -- (locality);
\draw[arrow] (locality) -- (gluing);
\draw[arrow] (gluing) -- (arithmetic);
\end{tikzpicture}

 \caption{The geometric proof of arithmetic exhaustion. Each local branch
 enters product calibration with its own speed and domain statement;
 simultaneous commutation is established only afterwards.}
 \label{fig:proof-map}
\end{figure}

\subsection{Organization}
Section~\ref{sec:models} verifies the elementary arithmetic models.
Sections~\ref{sec:setup}--\ref{sec:transfer} construct the finite-module
framework and ordered sharp boundary links. The finite-place, isolated
real, and higher-component real arguments occupy
Sections~\ref{sec:finite}--\ref{sec:panels}; Table~\ref{tab:local-branches}
records their exact inputs to Section~\ref{sec:heights}.
Sections~\ref{sec:locality}--\ref{sec:gluing} prove locality and
simultaneous commutation. Section~\ref{sec:recovery} extracts the
arithmetic models and completes Theorem~\ref{thm:arithmetic-exhaustion}.
All geometric and analytic arguments needed for this theorem are given
in this article.
